\documentclass[12pt]{article}
\usepackage{amsmath,amssymb}

\begin{document}
\section*{{2018 AIME II Problems/Problem 12}}
Source: AoPS Wiki

\subsection*{Solution}
Let $AP=x$ and let $PC=\rho x$ . Let $[ABP]=\Delta$ and let $[ADP]=\Lambda$ . We easily get $[PBC]=\rho \Delta$ and $[PCD]=\rho\Lambda$ . We are given that $[ABP] +[PCD] = [PBC]+[ADP]$ , which we can now write as \[\Delta + \rho\Lambda = \rho\Delta + \Lambda \qquad \Longrightarrow \qquad \Delta -\Lambda = \rho (\Delta -\Lambda).\] Either $\Delta = \Lambda$ or $\rho=1$ . The former would imply that $ABCD$ is a parallelogram, which it isn't; therefore we conclude $\rho=1$ and $P$ is the midpoint of $AC$ . Let $\angle BAD = \theta$ and $\angle BCD = \phi$ . Then $[ABCD]=2\cdot [BCD]=140\sin\phi$ . On one hand, since $[ABD]=[BCD]$ , we have \begin{align}\sqrt{65}\sin\theta = 7\sin\phi \quad \implies \quad 16+49\cos^2\phi = 65\cos^2\theta\end{align} whereas, on the other hand, using cosine formula to get the length of $BD$ , we get \[10^2+4\cdot 65 - 40\sqrt{65}\cos\theta = 10^2+14^2-280\cos\phi\] \begin{align}\tag{2}\implies \qquad 65\cos^2\theta = \left(7\cos\phi+ \frac{8}{5}\right)^2\end{align} Eliminating $\cos\theta$ in the above two equations and solving for $\cos\phi$ we get \[\cos\phi = \frac{3}{5}\qquad \implies \qquad \sin\phi = \frac{4}{5}\] which finally yields $[ABCD]=2\cdot [BCD] = 140\sin\phi = 112$ .

\end{document}
